\honest{The Quotation is not the Referent}{Eliezer Yudkowsky}{2008}

In logic, if A equals B, you may put one in place of the other: from (2 + 2) = 4 and ((2 + 2) + 3) = 7 follows (4 + 3) = 7. Apply this to beliefs and it goes wrong. The morning star and the evening star are both Venus. Mary believes the morning star is the god Lucifer. She thinks the evening star is the god Venus. John knows the two stars are one, and believes that Mary believes the morning star is Lucifer. Must he, by substitution, believe that Mary believes the evening star is Lucifer? The problem is ``usually phrased'' this way. \nb{The example and the puzzle are Frege's, from ``On Sense and Reference'' (1892). The post does not name him.}

A simpler version: ``2 + 2 = 4'' equals TRUE, Fermat's Last Theorem equals TRUE, so believing one means believing the other. I ask you to imagine someone writing a reasoning program that makes this mistake, writing a paper on how to prevent it, and someone else disagreeing. ``The argument is still going on.'' \nb{It was real. John McCarthy wrote in 1979 that ``The problem identified by Frege'' arises ``in artificial intelligence as well as in philosophy and linguistics''.}

``P'rsnally,'' I think John is committing a type error, like subtracting 5 grams from 20 meters; ``the morning star'' is not even the same type of thing as the morning star. ``Beliefs are not planets.'' The morning star equals the evening star, but ``morning star'' does not equal ``evening star''. The original mistake was an AI that stored its beliefs about Mary's beliefs about ``the morning star'' in the same form as its beliefs about the morning star. ``The whole paradox stems from the failure to use quote marks in appropriate places.'' \nb{Frege's own answer was close to this: after ``believes'', a term stands for a way of conceiving its object, not for the object. McCarthy's 1979 remedy was the same type distinction, and he described it as working ``without modal operators or quotation marks and without the restrictions on substituting equals for equals that either device makes necessary.'' Quotation was one of the rival proposals in the argument the post mentions; the post does not say which proposals it rejects, or why.}

I have urged type discipline before, in keeping expected utilities apart from utilities. It is like keeping track of units in physics, which seems a bother until your answer is off by orders of magnitude and comes out in seconds per square gram. Beliefs live in brains and weigh micrograms; planets live in space and weigh rather more. Since quote marks look too much like the words inside them, I suggest writing each letter as its number in the alphabet, so that ``morning star'' becomes 13.15.18.14.9.14.7.0.19.20.1.18.

In mathematical logic the statement that P is a theorem and the statement that it is provable that an encoded proof of the encoded sentence P exists are ``very distinct propositions.'' Dropping a level of quotation is like dropping a unit: you get results like ``The speed of light is 299,792,458 meters long.'' Tarski ``once tried to define the meaning of `true''' with an infinite family of sentences such as ``Snow is white'' is true if and only if snow is white, and ``Weasels are green'' is true if and only if weasels are green. \nb{The family was Tarski's test for a definition of truth; the definition was a single formula from which the sentences follow.} When such sentences start to seem meaningful, you have learned to tell sentences from the world.

Last, truth is not reality. ``Saying `true' compares a belief to reality,'' and reality need not be compared with any belief in order to be real. Remember this ``the next time someone claims that nothing is true.''
